\section{Sign clusters and exponential decay}\label{sec:clusters}

We now turn the uniform annulus estimate into a bound on connections.
The argument uses sign clusters, of the type studied for continuous-spin
systems in~\cite{CC}.  We prove the finite-volume comparison statements
needed here, including the comparison with pinned spins.

We return to the free spin law \eqref{eq:model} and its expectation
$\langle\cdot\rangle_{G,b}^{(n)}$, now on an arbitrary finite graph
$G=(V,E)$ with strengths $b_e\in[0,\beta]$.
Until the final subsection, $n=3$ and we write the spins as $s_x$.

\subsection{An associated amplitude and bond law}

Write
\begin{equation}\label{eq:spin-disintegration}
 s_x=(r_x\tau_x,q_x\cos\theta_x,q_x\sin\theta_x),
 \qquad q_x=(1-r_x^2)^{1/2}.
\end{equation}
Under the reference measure, the variables $r_x$, $\tau_x$, and
$\theta_x$ are independent and uniform on $[0,1]$, $\{-1,1\}$, and
$\R/(2\pi\Z)$, respectively.  Indeed, the first coordinate of a uniform
point on $S^2$ is uniform on $[-1,1]$, independently of its azimuthal
angle.  Conditional on all the amplitudes $r=(r_x)_{x\in V}$, the signs
and angles are independent Gibbs systems with respective couplings
\begin{equation}\label{eq:amplitude-couplings}
 K_{xy}=b_{xy}r_xr_y,
 \qquad L_{xy}=b_{xy}q_xq_y.
\end{equation}
Their edge energies are $\tau_x\tau_y$ and
$\cos(\theta_x-\theta_y)$.  Let $Z_I(K)$ and $Z_P(L)$ denote their
partition functions, using uniform probability reference measures.
The amplitude density is therefore proportional to
\begin{equation}\label{eq:amplitude-density}
 \rho(r)=Z_I(K(r))Z_P(L(r)).
\end{equation}

We now use the embedded Ising representation underlying Wolff's
reflection clusters~\cite{Wolff}, coupled to bonds as in
Edwards--Sokal~\cite{EdwardsSokal}. Given the spins, sample bonds
$\eta_e\in\{0,1\}$ independently. For $e=\{x,y\}$, set
\begin{equation}\label{eq:bond-parameter}
 p_e=1-e^{-2K_e},
 \qquad
 \mathbb P(\eta_e=1\mid s)
   =p_e\mathbf1_{\{\tau_x=\tau_y\}}.
\end{equation}
The elementary identity
\begin{equation}\label{eq:ising-bond-expansion}
 e^{K_e\tau_x\tau_y}
  =e^{K_e}\big[(1-p_e)+p_e\mathbf1_{\{\tau_x=\tau_y\}}\big]
\end{equation}
shows that, conditional on $r,\eta$, the signs are independent fair
signs on the open connected components. In particular, we obtain the
weighted connection identity of Campbell and Chayes~\cite[Corollary~4,
Equation~(7)]{CC}:
\begin{equation}\label{eq:spin-connection}
 \langle s_x^1s_y^1\rangle_{G,b}^{(3)}
 =\E\big[r_xr_y\mathbf1_{\{x\leftrightarrow y\}}\big],
\end{equation}
where the connection uses only open edges of $G$.

To bound these connections, we will show that fixing any set of spins
at $(1,0,0)$ cannot decrease the probability of an increasing bond
event. We will apply this comparison to simultaneous annulus crossings.
Its proof uses association jointly in amplitudes and bonds, because
the pinning condition constrains amplitudes as well as cluster signs.

We use a finite-product version of the
Fortuin--Kasteleyn--Ginibre association criterion~\cite{FKG}, allowing
both discrete and continuous coordinates, and include its proof.
A probability law on a coordinatewise ordered product is
\emph{associated} if the covariance of any two bounded increasing
functions is nonnegative.  A positive function $w$ is
\emph{log-supermodular} if
\[
 w(u\vee v)w(u\wedge v)\geq w(u)w(v),
\]
where join and meet are taken coordinatewise.

\begin{lemma}\label{lem:association}
Let $w$ be a positive log-supermodular density on a finite product of
compact intervals or finite chains, with respect to the product of
Lebesgue or counting measures.  In the interval case, suppose $w$ is
continuous.  The probability law proportional to $w$ is associated.
Moreover, its conditional law on all but one coordinate increases
stochastically as the remaining coordinate increases.
\end{lemma}

\begin{proof}
For one coordinate, association follows by taking independent copies
$X,X'$ and writing
\[
 2\operatorname{Cov}(F(X),H(X))
 =\E\big[(F(X)-F(X'))(H(X)-H(X'))\big]\geq0.
\]
Proceed by induction on the number of coordinates.  Write a point as
$(u,t)$, with $t$ the final coordinate.  Each conditional density
$w(u,t)$ is log-supermodular in $u$, and its law is associated by
induction.  For $t'\geq t$, log-supermodularity implies that
$w(u,t')/w(u,t)$ is increasing in $u$: apply the defining inequality to
$(u,t')$ and $(v,t)$ when $u\leq v$.  Tilting the conditional law at $t$
by this increasing likelihood ratio, and applying its association,
shows that the conditional law at $t'$ stochastically dominates it.
The ratios are bounded in the continuous case by positivity and
compactness.

If $F(u,t)$ and $H(u,t)$ are increasing, their conditional expectations
are consequently increasing functions of $t$.  The decomposition
\[
 \operatorname{Cov}(F,H)
 =\E\big[\operatorname{Cov}(F,H\mid t)\big]
  +\operatorname{Cov}\big(\E[F\mid t],\E[H\mid t]\big)
\]
has nonnegative terms by induction and the one-coordinate case.
This proves both assertions.
\end{proof}

The following restricted correlation inequalities are the Ising and
plane-rotator cases of the method of Ginibre~\cite{Ginibre}. Their
scope is important: no such inequality for the full vector-spin law
is assumed.

\begin{lemma}\label{lem:edge-correlation-inequalities}
In either the free Ising or the free planar system on a finite graph
with nonnegative couplings, every edge energy has nonnegative mean,
and any two edge energies have nonnegative covariance.
\end{lemma}

\begin{proof}
The mean inequality follows by expanding the Gibbs exponential.
For Ising spins, every reference moment of a sign monomial is either
zero or one.  For planar spins, each cosine of an angle difference has
nonnegative Fourier coefficients, so every product of such cosines
has a nonnegative constant Fourier coefficient.

For covariances, take two independent copies of the Gibbs system.
If $T_e$ and $T_f$ are two edge energies and primes denote the second
copy, then
\begin{equation}\label{eq:doubled-covariance}
 2\operatorname{Cov}(T_e,T_f)
 =\E\big[(T_e-T'_e)(T_f-T'_f)\big].
\end{equation}
In the planar reference integral, make the substitution
$\theta=\varphi+\psi$, $\theta'=\varphi-\psi$ modulo $2\pi$.
Independent uniform $\varphi,\psi$ produce independent uniform
$\theta,\theta'$: this is a surjective homomorphism of compact tori,
and therefore preserves Haar probability.  For an oriented edge $g$,
write $\varphi_g=\varphi_x-\varphi_y$, and similarly for $\psi_g$.
The doubled density numerator becomes
\[
 \exp\left(2\sum_{g\in E}L_g\cos\varphi_g\cos\psi_g\right),
 \qquad T_e-T'_e=-2\sin\varphi_e\sin\psi_e.
\]
Consequently the unnormalized numerator on the right of
\eqref{eq:doubled-covariance} is
\[
 4\sum_{a\in\N_0^E}
 \left(\prod_{g\in E}\frac{(2L_g)^{a_g}}{a_g!}\right)
 \left(
  \int\sin\varphi_e\sin\varphi_f
       \prod_{g\in E}(\cos\varphi_g)^{a_g}\,\dd\varphi
 \right)^2\geq0,
\]
where the integral uses product Haar probability.  Each integral is
real; the possible multiplicity of the torus map introduces no factor
because the measures are probabilities.

For Ising reference signs $d,d'$, set
$u=(d+d')/2$ and $v=(d-d')/2$ at each site.  Their sitewise values are
$(1,0),(-1,0),(0,1),(0,-1)$, with equal probabilities.  Every mixed
monomial in $u,v$ has nonnegative expectation: a monomial involving
positive powers of both vanishes, an odd pure moment vanishes, and
an even pure moment is nonnegative.  The doubled density is
\[
 \exp\left(2\sum_{g=\{x,y\}}K_g(u_xu_y+v_xv_y)\right),
\]
and an edge-energy difference is $2(u_xv_y+v_xu_y)$.
Expanding the exponential and the product of two differences gives
only nonnegative coefficients multiplying nonnegative reference
moments.  All expansions above are absolutely integrable on the finite
products, proving the covariance assertion.
\end{proof}

The next result is the finite-graph amplitude association established
by Campbell and Chayes~\cite[Lemma~2]{CC}. We include its proof because
both its sign structure and its boundary-free hypothesis matter below.

\begin{lemma}\label{lem:amplitude-association}
The amplitude law with density \eqref{eq:amplitude-density} is associated.
\end{lemma}

\begin{proof}
For either partition function, with couplings $J_e(r)$ and edge
energies $T_e$, differentiation in distinct site variables $r_x,r_y$
gives
\begin{align*}
 \partial_x\partial_y\log Z(J(r))
 &=\sum_e\langle T_e\rangle\,\partial_x\partial_y J_e\\
 &\quad+\sum_{e,f}\operatorname{Cov}(T_e,T_f)
              (\partial_xJ_e)(\partial_yJ_f).
\end{align*}
For $J=K$, the first derivatives are nonnegative.  For $J=L$, they are
nonpositive, so the products in the second sum are again nonnegative.
The distinct-site mixed derivatives of individual couplings vanish
unless $e=\{x,y\}$, in which case
\[
 \partial_x\partial_y K_{xy}=b_{xy},
 \qquad
 \partial_x\partial_y L_{xy}
   =b_{xy}\frac{r_xr_y}{q_xq_y}\geq0
\]
in the open amplitude cube.  Lemma~\ref{lem:edge-correlation-inequalities}
therefore shows that every distinct-site mixed derivative of
$\log\rho$ is nonnegative there.  An increment of $\log\rho$ in one
coordinate is thus nondecreasing in every other coordinate.  For two
points $u,v$, compare the successive increments from $u\wedge v$ to
$u$ with the same coordinate increments from $v$ to $u\vee v$.
Their sum gives the log-supermodular inequality in the open cube,
and continuity extends it to the closed cube.  The possible singularity
of derivatives at $r_x=1$ is immaterial.  The density itself is positive
and continuous, so Lemma~\ref{lem:association} applies.
\end{proof}

The conditioning argument used for bond association in
Campbell--Chayes~\cite[Corollary~1]{CC} also gives the following joint
statement.

\begin{lemma}\label{lem:joint-association}
The joint free law of $r$ and $\eta$ is associated, with both amplitudes
and bond indicators ordered coordinatewise.
\end{lemma}

\begin{proof}
Conditional on $r$, summing the signs in
\eqref{eq:ising-bond-expansion} gives a bond weight proportional to
\begin{equation}\label{eq:conditional-fk-weight}
 2^{k(\eta)}\prod_{e\in E}p_e^{\eta_e}(1-p_e)^{1-\eta_e},
\end{equation}
where $k(\eta)$ counts all open components, including singletons.
The function $k$ is supermodular: the change on opening an edge is
$-1$ if its endpoints were disconnected and $0$ otherwise, and this
change increases as other bonds are opened.  Hence the weight is
log-supermodular and its law is associated by
Lemma~\ref{lem:association}.  Increasing any nondegenerate parameter
$p_e$ tilts the weight by an increasing likelihood ratio.  Association
therefore implies stochastic increase of the conditional bond law
in each $p_e$, and consequently in $r$.  A zero parameter forces its
edge closed; restriction to the remaining edges, or continuity, handles
these degenerate cases.

For bounded increasing functions $F(r,\eta)$ and $H(r,\eta)$,
conditional association makes
$\E[\operatorname{Cov}(F,H\mid r)]$ nonnegative.  Their conditional
means are increasing in $r$, by stochastic monotonicity of the bond
law and their direct monotonicity in $r$.  Their covariance is
nonnegative by Lemma~\ref{lem:amplitude-association}.  The total
covariance identity proves the claim.
\end{proof}

\subsection{Pinning and crossings}

Let $s_*=(1,0,0)$.  For any set $B\subseteq V$, the pinned spin law
means that $s_x=s_*$ for $x\in B$, while all other spins retain uniform
reference measures and the same edge interactions.  Bonds are sampled
by the same rule \eqref{eq:bond-parameter}.

\begin{proposition}[Pinning comparison]\label{prop:pinning}
For the three-component free spin law on any finite graph with
nonnegative strengths, and bonds sampled by \eqref{eq:bond-parameter},
pinning any set of spins to $s_*$ cannot decrease the probability of
any increasing bond event $A$.
\end{proposition}

\begin{proof}
For $0<\varepsilon<1$, condition the free system on
$r_x\geq1-\varepsilon$ and $\tau_x=1$ for every $x\in B$.
Conditional on $r,\eta$, the probability of the sign requirement is
$2^{-k_B(\eta)}$, where $k_B$ is the number of open components meeting
$B$.  Thus the induced law of $r,\eta$ is tilted by
\[
 T_\varepsilon(r,\eta)
 =\mathbf1_{\{r_x\geq1-\varepsilon\ \text{for all }x\in B\}}
       2^{-k_B(\eta)}.
\]
Opening an edge cannot increase $k_B$; this also holds when it merges
a component meeting $B$ with one not meeting $B$.  Consequently
$T_\varepsilon$ is increasing.  Lemma~\ref{lem:joint-association} gives
\[
 \frac{\E[\mathbf1_A T_\varepsilon]}{\E[T_\varepsilon]}
 \geq\mathbb P(A).
\]

As $\varepsilon\downarrow0$, the conditioned reference measures on
the positive caps converge to point masses at $s_*$.  The finite-volume
Gibbs weight is bounded, positive, and continuous.  The spin-to-bond
kernel is also continuous, since its edge-opening probability is
\[
 1-\exp\big(-2b_{xy}\max\{s_x^1s_y^1,0\}\big).
\]
In particular there is no discontinuity at a zero first coordinate.
Integrating any event of the finitely many bonds and passing to the
limit proves the asserted comparison for exact pinning.
\end{proof}

We now specialize to finite subgraphs of the nearest-neighbor graph
on $\Z^2$, with arbitrary strengths in $[0,\beta]$.  For $z\in\Z^2$
and an integer $N\geq2$, put
\[
 D_N(z)=\{v\in\Z^2:N\leq\|v-z\|_\infty\leq2N\}.
\]
An \emph{annulus crossing} is an open path in the restriction of $G$
to $D_N(z)$, joining its layers of radii $N$ and $2N$.
Edges or vertices outside this annulus are not allowed in the path.

\begin{proposition}[Uniform crossing suppression]\label{prop:crossings}
For each $\beta<\infty$ there are numbers $p_N=p_N(\beta)>0$ tending
to zero as $N\to\infty$ with the following property.  For every finite
lattice subgraph, every strength array in $[0,\beta]$, and every list
of centers whose closed boxes of sup radius $2N$ are pairwise disjoint,
the free probability that all the corresponding annuli are crossed
is at most $p_N$ to the number of annuli.  The boxes need not lie
inside the finite graph.
\end{proposition}

\begin{proof}
First work in one restricted annulus graph.  Pin both of its layers
to $s_*$, and let $Z(0)$ be its partition function.  Let $Z(\pi)$ be
the partition function with the inner layer changed to $-s_*$ and
the outer layer left at $s_*$.  Then
\begin{equation}\label{eq:crossing-partition-ratio}
 \mathbb P_{\mathrm{pinned}}(\text{crossing})
   =1-\frac{Z(\pi)}{Z(0)}.
\end{equation}
Indeed, in the expansion \eqref{eq:ising-bond-expansion}, every open
component missing the pinned layers has two sign choices.  A component
meeting a pinned layer has one sign choice, provided all prescribed
signs on it agree.  Changing the inner signs therefore removes exactly
the bond configurations joining the two layers and changes the weight
of none of the remaining configurations.  All other factors are
unchanged: the boundary amplitudes are $r=1,q=0$, and the reference
normalization for the unpinned signs is the same.  This argument holds
at every interior amplitude configuration and hence after integration.
If a layer has no vertices, both sides of
\eqref{eq:crossing-partition-ratio} are zero.

By Proposition~\ref{prop:annulus}, the right side of
\eqref{eq:crossing-partition-ratio} is bounded uniformly by a
nonnegative sequence tending to zero.  Enlarge it slightly to obtain
a strictly positive sequence $p_N$ with the same limit.

For several annuli, pin all their present boundary-layer sites to
$s_*$ simultaneously.  The intersection of their crossing events is
increasing, so Proposition~\ref{prop:pinning} bounds its free
probability by its pinned probability.  In the pinned system, each
unpinned annulus region is separated from every vertex outside that
annulus by the pinned layers: a nearest-neighbor path cannot jump
across either layer.  The unpinned annulus spins are therefore
independent with precisely the required fixed boundary values.
For each resulting annulus marginal, apply
Equation~\eqref{eq:crossing-partition-ratio} to the partition functions
of that restricted annulus graph.
Boundary-only interactions contribute constants, and the independent
sampling of edge bonds preserves independence of the crossing events.
Their probabilities multiply, proving the bound.

This separation uses only edges present in $G$.  Missing vertices or
edges, including those beyond the finite volume, cannot create an
unseparated nearest-neighbor interaction.  Pinning the layer sites
that are present and applying Proposition~\ref{prop:annulus} to the
restricted subgraph proves the assertion also for truncated boxes.
\end{proof}

\subsection{From crossings to exponential decay}

\begin{proposition}\label{prop:connection}
For every finite $\beta>0$ there are constants $A_0<\infty$ and
$m_0>0$, depending only on $\beta$, such that the three-component free
bond law on every finite square-lattice subgraph with strengths in
$[0,\beta]$ satisfies
\[
 \mathbb P(x\leftrightarrow y)
 \leq A_0\exp(-m_0\|x-y\|_2)
 \qquad(x,y\in V).
\]
\end{proposition}

\begin{proof}
Translate the lattice so that $x=0$.  Choose a finite integer $N\geq2$
so large that, for the sequence in Proposition~\ref{prop:crossings},
\begin{equation}\label{eq:coarse-path-alpha}
 0<p:=p_N<1,
 \qquad \alpha:=4p^{1/25}<1.
\end{equation}
Partition $\Z^2$ into boxes
$z+\{0,\ldots,N-1\}^2$, with coarse vertices $z\in N\Z^2$.
Suppose $0$ and $y$ are joined by an open path and
$d:=\|y\|_\infty>4N$.  Recording the visited boxes and erasing loops
produces a simple nearest-neighbor coarse path from the origin's box
to the box containing $y$.  If its length is $k$, then
\begin{equation}\label{eq:coarse-path-length}
 k\geq d/N-1.
\end{equation}

Each vertex $z$ of this simple coarse path corresponds to a box
visited by the original open path, at some vertex strictly within
sup distance $N$ of $z$.  At least one of the original endpoints is
at distance greater than $2N$ from $z$: otherwise their mutual
distance would be at most $4N$.  A segment of the original path from
the visited vertex to that endpoint thus contains a crossing of
$D_N(z)$: take its first hit on the outer layer and its last preceding
hit on the inner layer. The intervening segment stays in the annulus.  It does not matter whether that segment survived coarse
loop erasure; its edges are still open.

Among the $k+1$ distinct coarse vertices, one of the 25 residue classes
modulo $5N$ in the two coordinates contains at least $(k+1)/25$
vertices.  Their closed outer boxes of sup radius $2N$ are disjoint,
because distinct centers in the class differ by at least $5N$ in
some coordinate.  Choose one maximizing class deterministically for
each candidate coarse path.  Proposition~\ref{prop:crossings} bounds
the probability that this path has all its required crossings by
$p^{(k+1)/25}$.  There are at most $4^k$ coarse paths of length $k$
from the initial box.  The union bound and
\eqref{eq:coarse-path-length} therefore give
\begin{align*}
 \mathbb P(0\leftrightarrow y)
 &\leq p^{1/25}\sum_{k\geq\lceil d/N-1\rceil}\alpha^k\\
 &\leq \frac{p^{1/25}}{\alpha(1-\alpha)}
           \exp\left(-\frac{-\log\alpha}{N}\,d\right).
\end{align*}
Since $d\geq\|y\|_2/\sqrt2$, take
\[
 m_0=\frac{-\log\alpha}{\sqrt2\,N}>0.
\]
Increasing the finite prefactor to cover $d\leq4N$ proves the result.
Every estimate was uniform in the finite subgraph and its strengths.
\end{proof}

\subsection{Conditioning to three components}

We have proved the uniform connection estimate in three dimensions.
To pass to higher-dimensional spheres, we need a statement about the
conditional reference measure, not a new association theorem.
The following is the block-coordinate version of the conditional
projection in Aru--Garban--Sep\'ulveda~\cite[Proposition~2.4]{AGS}.

\begin{lemma}[Three-component conditional law]\label{lem:conditioning}
Let $n>3$, and let $\sigma_x\in S^{n-1}$ have the free Gibbs law
\eqref{eq:model} on a finite graph. Write
$\sigma_x=(\rho_x s_x,v_x)$, with $v_x\in\R^{n-3}$,
$\rho_x=\sqrt{1-|v_x|^2}$, and $s_x\in S^2$ whenever $\rho_x>0$.
Conditional on $(v_x)_{x\in V}$, the directions $(s_x)_{x\in V}$ have
exactly the free three-component Gibbs law with strengths
$b_{xy}\rho_x\rho_y\in[0,\beta]$.
\end{lemma}

\begin{proof}
Represent a uniform sphere point by
$(U,V)/\sqrt{|U|^2+|V|^2}$, where $U\in\R^3$ and
$V\in\R^{n-3}$ have independent standard Gaussian coordinates.
The direction $U/|U|$ is uniform on $S^2$ and is independent of
$(|U|,V)$. Since $v=V/\sqrt{|U|^2+|V|^2}$ depends only on these
latter variables, the three-component direction is independent of
$v$ under the reference law. Applying this representation independently
at all sites proves that, conditional on all $v_x$, the reference
directions remain independent and uniform on $S^2$.

The conditioned interaction is
\[
 \sum_{\{x,y\}\in E}b_{xy}v_x\cdot v_y
 +\sum_{\{x,y\}\in E}b_{xy}\rho_x\rho_y s_x\cdot s_y.
\]
The first sum is constant under the conditional law and cancels from
its normalization. The second gives the asserted couplings. The event
$\rho_x=0$ has reference probability zero and still has Gibbs
probability zero because the finite Gibbs density is positive and
bounded. Alternatively, choosing any independent reference direction
at such a site would decouple it and leave all spin products unchanged.
\end{proof}

\begin{proof}[Proof of Theorem~\ref{thm:main} and Corollary~\ref{cor:limits}]
For $n=3$, Equation~\eqref{eq:spin-connection} and
Proposition~\ref{prop:connection} give
\[
 0\leq\langle s_x^1s_y^1\rangle_{G,b}^{(3)}
 \leq A_0e^{-m_0\|x-y\|_2}.
\]
For $n>3$, apply Lemma~\ref{lem:conditioning}. Conditional on all
$v_x$, the first-coordinate correlation is $\rho_x\rho_y$ times
the corresponding three-component first-coordinate correlation.
The latter is nonnegative and satisfies the same bound, uniformly
in the conditional strength array. Since $0\leq\rho_x\rho_y\leq1$,
averaging gives, for every $n\ge3$,
\[
 0\leq\langle\sigma_x^1\sigma_y^1\rangle_{G,b}^{(n)}
 \leq A_0e^{-m_0\|x-y\|_2}.
\]
The original free spin law is invariant under coordinate permutations,
so
\[
 \langle\sigma_x\cdot\sigma_y\rangle_{G,b}^{(n)}
 =n\langle\sigma_x^1\sigma_y^1\rangle_{G,b}^{(n)}\geq0.
\]
Taking $A=nA_0$ and $m=m_0$ proves Theorem~\ref{thm:main}.

The constants are independent of the volume, so the same estimate
passes to the limsup in~\eqref{eq:limsup}. If a subsequence converges
locally weakly, the expectation of the bounded continuous local
function $\sigma\mapsto\sigma_0\cdot\sigma_x$ converges along it.
Its limit therefore satisfies the same bound. This proves
Corollary~\ref{cor:limits} without a uniqueness assertion.
\end{proof}
